\chapter{Week 6: Mutable State}
\label{ch:week6}

Lecture 6: October 13, 2026. Recitation~6 will be added after it meets.

Supplementary reading:
\begin{itemize}
  \item \cite{PLAI} Chapter 8 (Mutation: Structures and Variables), Section 8.1 in particular, which adds boxes to the closure interpreter with a store as this lecture does.
  \item \cite{SICP} Section 3.1 (Assignment and Local State), for what assignment costs: an expression can no longer be replaced by its value.
  \item \cite{Reynolds}, optional: the paper that names the problem of \Cref{sec:metacirc}, an interpreter that silently inherits features of the language it is written in.
\end{itemize}

Every interpreter so far has been \emph{pure}: evaluating an expression produced a value and nothing else, so the same expression in the same environment always gave the same answer. Lecture~6 adds the smallest piece of mutable state, the box of \Cref{sec:scopestate}, to last week's closure interpreter (\code{lecture5e.rkt}, \Cref{sec:closures}). The first attempt stores a box's contents directly in the value and cannot express \code{set-box!}. The fix is a second table next to the environment, the \emph{store}, which maps locations to values and is passed into and out of every evaluation step. The lecture then writes the same interpreter a second time using Racket's own boxes, and compares the two. This chapter covers:
\begin{itemize}
  \item \Cref{sec:stateboxes}: the four new constructs, the examples, and a first attempt that cannot mutate;
  \item \Cref{sec:storepass}: store-passing style, with locations, \code{fetch}, the \tyname{Result} type, and threading the store through arithmetic, \code{let}, and application;
  \item \Cref{sec:storeboxes}: the box clauses of the store-passing interpreter, with traces showing where the environment and where the store come from;
  \item \Cref{sec:metacirc}: the meta-circular interpreter, which implements boxes with Racket boxes, and what each version explains.
\end{itemize}

\section{State and Boxes}
\label{sec:stateboxes}

\subsection{Four New Constructs}

The starting file is ``a starting point for today's exploration of mutable state''. It takes the closure interpreter of \code{lecture5e.rkt} (\Cref{sec:closures}) and adds four constructs to \tyname{Expr}: \code{boxC} creates a box, \code{unboxC} reads its contents, \code{setboxC} updates them, and \code{seqC} is Racket's \code{(begin e1 e2)}.
\begin{lstlisting}
(define-type Expr
  [numC (n : number)]
  [plusC (e1 : Expr) (e2 : Expr)]
  [timesC (e1 : Expr) (e2 : Expr)]
  [letC (x : symbol) (e1 : Expr) (e2 : Expr)]
  [lambdaC (x : symbol) (e : Expr)]
  [appC (e1 : Expr) (e2 : Expr)]
  [idC (x : symbol)]
  [boxC (e : Expr)]
  [unboxC (e : Expr)]
  [setboxC (e1 : Expr) (e2 : Expr)]
  [seqC (e1 : Expr) (e2 : Expr)])
\end{lstlisting}
The parser gets one line per construct, in the \code{case} on the head symbol, with Racket's names as the concrete syntax:
\begin{lstlisting}
            [(begin) (seqC (parse (second l)) (parse (third l)))]
            [(box) (boxC (parse (second l)))]
            [(set-box!) (setboxC (parse (second l)) (parse (third l)))]
            [(unbox) (unboxC (parse (second l)))]
\end{lstlisting}
\begin{center}
\begin{tabular}{@{}lll@{}}
\toprule
concrete syntax & abstract syntax & meaning \\
\midrule
\code{(box e)} & \code{boxC} & a new box whose contents is the value of \code{e} \\
\code{(unbox e)} & \code{unboxC} & the current contents of the box \code{e} \\
\code{(set-box! e1 e2)} & \code{setboxC} & make the value of \code{e2} the contents of the box \code{e1} \\
\code{(begin e1 e2)} & \code{seqC} & evaluate \code{e1}, then \code{e2}; the value is that of \code{e2} \\
\bottomrule
\end{tabular}
\end{center}
\code{begin} takes exactly two expressions, so a longer sequence is written with nested \code{begin}s. That is all the starting file provides: \tyname{Value} is still \code{numV} and \code{closV}, and \code{eval-env} is the one of \code{lecture5e.rkt} with an \code{else} clause that reports the new variants as unimplemented. The rest of the lecture writes the four clauses, twice.

\subsection{The Examples}

The file ends with five test programs:
\begin{lstlisting}
(define ex1 (parse '(begin (+ 1 1) (+ 2 2))))
(define ex2 (parse '(let b (box 0) (set-box! b 1))))
(define ex3 (parse '(let b (box 0) (unbox b))))
(define ex4 (parse '(let b (box 0) (begin (begin (set-box! b 1)
                                                 (set-box! b (+ 1 (unbox b))))
                                          (unbox b)))))

(define ex5
  (parse '(let a (box 1)
            (let f (lambda x (+ x (unbox a)))
              (begin
                (set-box! a 2)
                (f 10))))))
\end{lstlisting}
Read as Racket programs, they should give:
\begin{itemize}
  \item \code{ex1}: 4. The value of \code{(+ 1 1)} is computed and thrown away.
  \item \code{ex2}: whatever \code{set-box!} returns. Racket's returns \code{(void)}; this language will return the box (\Cref{sec:storeboxes}).
  \item \code{ex3}: 0, the initial contents.
  \item \code{ex4}: 2. The box is set to 1, then to one more than its contents, and read.
  \item \code{ex5}: 12. This is the interesting one. The closure \code{f} is created while \code{a} holds 1, and called after \code{a} has been set to 2. The closure saved its environment (\Cref{sec:closures}), yet \code{(unbox a)} in its body must see the \emph{new} contents, so $10 + 2 = 12$. Racket itself gives 12 for the same program.
\end{itemize}

\subsection{Attempt 1: A Box Holds Its Value}

The obvious first try: a box value simply contains the value it holds.
\begin{lstlisting}
(define-type Value
  [numV (n : number)]
  [closV (env : Env) (x : symbol) (e : Expr)]
  [boxV (v : Value)])

    [boxC (e) (boxV (eval-env env e))]
    [unboxC (e) (boxV-v (eval-env env e))]
    [setboxC (e1 e2)
             (let ([v (eval-env env e1)])
               (boxV (eval-env env e2)))]
    [seqC (e1 e2)
          (let ([v (eval-env env e1)])
            (eval-env env e2))]
\end{lstlisting}
\code{boxC} and \code{unboxC} are fine: wrap the value, unwrap it. The trouble is \code{setboxC}. It evaluates the box \code{e1} into \code{v}, and then there is nothing to do with \code{v}: a \code{boxV} is an immutable piece of data, so the only option is to build a \emph{new} \code{boxV} holding the new value. \code{seqC} evaluates \code{e1}, ignores the result, and evaluates \code{e2}. On the examples:
\begin{center}
\begin{tabular}{@{}lll@{}}
\toprule
& Attempt 1 & expected \\
\midrule
\code{ex1} & \code{(numV 4)} & 4 \\
\code{ex2} & \code{(boxV (numV 1))} & the box, now holding 1 \\
\code{ex3} & \code{(numV 0)} & 0 \\
\code{ex4} & \textcolor{sigc}{\code{(numV 0)}} & 2 \\
\code{ex5} & \textcolor{sigc}{\code{(numV 11)}} & 12 \\
\bottomrule
\end{tabular}
\end{center}
In \code{ex4}, each \code{set-box!} builds a new box that the enclosing \code{begin} throws away, while the environment still binds \code{b} to the original \code{(boxV (numV 0))} (\Cref{fig:attempt1}). The final \code{(unbox b)} reads that box and gets 0. In \code{ex5}, \code{f}'s saved environment binds \code{a} to \code{(boxV (numV 1))} forever, so the call gives 11.

\begin{figure}[htbp]
\centering
\begin{tikzpicture}
  \node[cell, fill=greybg, minimum width=10mm] (b) at (0,0) {b};
  \node[lbl, above=1mm of b] {env};
  \node[cell, fill=tsk, minimum width=30mm] (v0) at (3.4,0) {(boxV (numV 0))};
  \draw[ptr] (b) -- (v0);
  \node[lbl, right=2mm of v0, align=left] {what every \code{(unbox b)} reads};
  \node[cell, fill=kspace, minimum width=30mm] (v1) at (3.4,-1.2) {(boxV (numV 1))};
  \node[lbl, right=2mm of v1, align=left] {built by \code{(set-box! b 1)},\\ then dropped by \code{begin}};
  \node[cell, fill=kspace, minimum width=30mm] (v2) at (3.4,-2.4) {(boxV (numV 1))};
  \node[lbl, right=2mm of v2, align=left] {built by \code{(set-box! b (+ 1 (unbox b)))},\\ also dropped};
\end{tikzpicture}
\caption{Attempt 1 on \code{ex4}. Nothing points to the boxes that \code{set-box!} builds; the name \code{b} still refers to the original, unchanged box. The second one holds 1 because \code{(unbox b)} still reads 0.}
\label{fig:attempt1}
\end{figure}

\begin{pitfall}[\code{ex2} looks right]
\code{(eval ex2)} is \code{(boxV (numV 1))}, exactly what a working \code{set-box!} would return, and \code{(unbox (set-box! (box 0) 7))} gives \code{(numV 7)}. The new contents are visible through the \emph{return value} of \code{set-box!}, never through the name of the box. That is functional update, like \code{cons} building a new list, not mutation. A test of mutation has to read the box again through another path afterwards, as \code{ex4} and \code{ex5} do.
\end{pitfall}

The same failure shows up whenever two parts of a program hold the same box. In \code{(let a (box 0) (let b a (begin (set-box! b 3) (unbox a))))}, \code{a} and \code{b} name one box, so the answer should be 3; Attempt 1 gives \code{(numV 0)}. A binding in the environment holds a \emph{copy} of a value, and values are immutable, so a box cannot be a value that contains its contents.

\begin{keypoint}[Mutation needs a level of indirection]
For a change made through one reference to be seen through another, both must refer to a shared place that holds the contents, rather than contain the contents themselves. The box value must say \emph{where} its contents are, and the contents must be kept somewhere that every part of the evaluation can see and that can change over time.
\end{keypoint}

\section{Store-Passing Style}
\label{sec:storepass}

\subsection{Locations and the Store}

The indirection is a second table. A box value no longer contains its contents; it contains a \emph{location}, a number that names a cell in a table called the \emph{store}.

\begin{definition}[Location, store]
A \emph{location} is an address, here a number. The \emph{store} maps locations to values. A box value is a location; its contents are whatever the store currently maps that location to.
\end{definition}
So a name now goes through up to two tables: the environment maps the name to a value, and if that value is a box, the store maps the box's location to its contents. The code for the store copies the code for the environment (\Cref{sec:env}) line for line:
\begin{lstlisting}
(define-type-alias Location number)

(define-type Storage
  [cell (l : Location) (val : Value)])

(define-type-alias Store (listof Storage))
(define empty-store empty)
(define override-store cons)

(define (fetch (l : Location) (sto : Store)) : Value
  (cond
    [(cons? sto)
     (if (equal? (cell-l (first sto)) l)
         (cell-val (first sto))
         (fetch l (rest sto)))]
    [else (error 'fetch "No location found")]))
\end{lstlisting}
\begin{center}
\begin{tabular}{@{}lll@{}}
\toprule
& environment & store \\
\midrule
maps & names (\code{symbol}) to values & locations (\code{number}) to values \\
entry & \code{(bind name val)} & \code{(cell l val)} \\
add an entry & \code{extend-env} = \code{cons} & \code{override-store} = \code{cons} \\
look up & \code{lookup}, newest first & \code{fetch}, newest first \\
\bottomrule
\end{tabular}
\end{center}
\code{fetch} is \code{lookup} with the names changed. Writing to a location does not remove its old cell: \code{override-store} puts a new cell for the same location on the front, and \code{fetch} stops at the first match. The latest write wins in the same way that the innermost binding wins in \code{lookup} (\Cref{fig:envstack}), with shadowing in time instead of shadowing in the program text.

\subsection{Results Carry the Store}

A box value is now a location:
\begin{lstlisting}
(define-type Value
  [numV (n : number)]
  [closV (env : Env) (x : symbol) (e : Expr)]
  [boxV (l : Location)])
\end{lstlisting}
Evaluating an expression may now change the store, and whoever evaluates next must see the change. So \code{eval-env} takes the current store as an argument and returns, together with the value, the store as it is \emph{after} the evaluation:
\begin{lstlisting}
(define-type Result
  [res (v : Value) (sto : Store)])

(define (eval-env (env : Env) (sto : Store) (e : Expr)) : Result
  (type-case Expr e
    [numC (n) (res (numV n) sto)]
    [lambdaC (x e) (res (closV env x e) sto)]
    [idC (x) (res (lookup x env) sto)]
    ...))

(define (eval (e : Expr))
  (eval-env empty-env empty-store e))
\end{lstlisting}
Numbers, \code{lambda}, and names do not touch the store, so they return the store they were given.

\begin{keypoint}[Store-passing style]
Every evaluation takes the store as input and returns the updated store as output. An effect is no longer something that \emph{happens}; it is a difference between the store that went in and the store that came out. The interpreter itself stays pure: no Racket mutation is needed to explain the language's mutation (with one exception, \code{new-loc} below).
\end{keypoint}

\subsection{Fresh Locations}

\code{(box e)} needs a location that no existing box uses. The lecture first wrote \code{new-loc} like this, and then commented it out:
\begin{lstlisting}
#|
(define new-loc
  (let ([b (box 0)])
    (let ([l (unbox b)])
      (begin
        (set-box! b (+ 1 l))
        l))))
|#
\end{lstlisting}
There is no \code{lambda}, so the right-hand side is evaluated once, when the \code{define} runs: it creates the box, reads 0, sets the box to 1, and returns 0. \code{new-loc} is just the number 0, and the box is never seen again. Asking for \code{new-loc} twice gives 0 both times, and \code{(new-loc)} does not even type check:
\begin{lstlisting}[style=repl]
> new-loc
- number
0
> (new-loc)
typecheck failed: call of a non-function
\end{lstlisting}
The fix wraps the body in \code{(lambda () ...)} inside the \code{let}:
\begin{lstlisting}
(define new-loc
  (let ([b (box 0)])
    (lambda ()
      (let ([l (unbox b)])
        (begin
          (set-box! b (+ 1 l))
          l)))))
\end{lstlisting}
Now the \code{let} runs once and creates one box, and \code{new-loc} is a closure over it; the body runs at every call. This is the counter pattern: \code{make-adder} of \Cref{sec:hof} with a box in place of a number, so the closure's private state can change.
\begin{lstlisting}[style=repl]
> (new-loc)
- number
0
> (new-loc)
- number
1
> (new-loc)
- number
2
\end{lstlisting}

\begin{remark}
\code{new-loc} implements the language's mutation using Racket's: \code{set-box!} on a Racket box. It is the one place where the store-passing interpreter is not pure. Only freshness matters, so the effect is harmless, but the counter is global: it is never reset between programs. Evaluating \code{ex2} twice in one session gives \code{(res (boxV 0) ...)} the first time and \code{(res (boxV 1) ...)} the second, with the same contents. The traces in this chapter are each run in a fresh session, so locations start at 0.
\end{remark}

\subsection{Threading the Store Through Arithmetic}

\code{plusC} has two subexpressions, either of which may change the store. The clause evaluates \code{e1} in the incoming store \code{sto}, giving \code{sto1}; evaluates \code{e2} in \code{sto1}, giving \code{sto2}; and returns \code{sto2}:
\begin{lstlisting}
    [plusC (e1 e2)
           (let ([r1 (eval-env env sto e1)])
             (type-case Result r1
               [res (v1 sto1)
                    (let ([r2 (eval-env env sto1 e2)])
                      (type-case Result r2
                        [res (v2 sto2)
                             (res (numV (+ (numV-n v1) (numV-n v2))) sto2)]))]))]
\end{lstlisting}
\code{type-case Result} is used here only to take a \code{res} apart into its two fields. \code{timesC} is the same with \code{*}. The store is \emph{threaded}: each evaluation's output store is the next one's input (\Cref{fig:thread}), while \code{env} is passed unchanged to both operands, as before.

\begin{figure}[htbp]
\centering
\begin{tikzpicture}[x=1cm, y=1cm]
  \node[cell, fill=greybg, minimum width=10mm] (s0) at (0,0) {sto};
  \node[box, fill=kspace, font=\ttfamily] (e1) at (2.4,0) {eval e1};
  \node[cell, fill=greybg, minimum width=10mm] (s1) at (4.8,0) {sto1};
  \node[box, fill=kspace, font=\ttfamily] (e2) at (7.2,0) {eval e2};
  \node[cell, fill=greybg, minimum width=10mm] (s2) at (9.6,0) {sto2};
  \draw[arr] (s0) -- (e1); \draw[arr] (e1) -- (s1);
  \draw[arr] (s1) -- (e2); \draw[arr] (e2) -- (s2);
  \node[cell, fill=uspace, minimum width=10mm] (env) at (4.8,1.5) {env};
  \draw[ptr, densely dashed] (env) -| (e1);
  \draw[ptr, densely dashed] (env) -| (e2);
  \node[lbl, below=1mm of e1] {value \code{v1}};
  \node[lbl, below=1mm of e2] {value \code{v2}};
  \node[lbl, below=1mm of s2, align=center] {returned with\\ \code{(numV (+ v1 v2))}};
\end{tikzpicture}
\caption{Threading the store through \code{(+ e1 e2)}. The store flows through the operands in sequence; the environment is handed to both unchanged.}
\label{fig:thread}
\end{figure}

\begin{pitfall}[Passing the wrong store]
Every store has type \tyname{Store}, so the type checker accepts any of \code{sto}, \code{sto1}, \code{sto2} in any position. Two plausible slips in \code{plusC} are evaluating \code{e2} in \code{sto}, and returning \code{sto1}. Each is caught by a program that writes to a box inside one operand (these variants are not lecture code):
\begin{lstlisting}
(let b (box 0) (+ (begin (set-box! b 5) 0) (unbox b)))   ; P1
(let b (box 0) (begin (+ 0 (begin (set-box! b 5) 0)) (unbox b)))   ; P2
\end{lstlisting}
\begin{center}
\begin{tabular}{@{}lccc@{}}
\toprule
program & lecture \code{plusC} & \code{e2} gets \code{sto} & returns \code{sto1} \\
\midrule
P1 & \code{(numV 5)} & \textcolor{sigc}{\code{(numV 0)}} & \code{(numV 5)} \\
P2 & \code{(numV 5)} & \code{(numV 5)} & \textcolor{sigc}{\code{(numV 0)}} \\
\bottomrule
\end{tabular}
\end{center}
Evaluating \code{e2} in \code{sto} makes \code{e2} blind to \code{e1}'s writes; returning \code{sto1} forgets \code{e2}'s writes once the addition is over. Each store must be used exactly once, by the next step.
\end{pitfall}

\subsection{\texttt{let} Is Application, Again}

\code{letC} would need the same threading: evaluate \code{e1} in \code{sto}, then the body in the resulting store. The lecture avoids writing it by turning \code{let} into application, the equation of \Cref{sec:substrev}:
\begin{lstlisting}
    [letC (x e1 e2) (eval-env env sto (appC (lambdaC x e2) e1))]
\end{lstlisting}
The rewriting happens in \code{eval-env} rather than in \code{parse}, but the effect is the same: the store is threaded through \code{let} by the code for \code{appC}, written once. The equation is valid here because both \code{let} and application are eager.

\subsection{Application}

\code{appC} evaluates the function in \code{sto} and the argument in \code{sto1}, as \code{plusC} did, and then runs the body:
\begin{lstlisting}
    [appC (e1 e2)
          (type-case Result (eval-env env sto e1)
            [res (v1 sto1)
                 (type-case Result (eval-env env sto1 e2)
                   [res (v2 sto2)
                        (eval-env (extend-env (bind (closV-x v1) v2) (closV-env v1))
                                  sto2
                                  (closV-e v1))])])]
\end{lstlisting}
The body gets its environment and its store from different places. The environment is the closure's, \code{(closV-env v1)}, extended with the parameter: lexical scope, unchanged from \code{lecture5e.rkt}, whose comment the clause keeps (``we're going to use and extend the saved environment from the closure, instead of the current environment''). The store is \code{sto2}, the \emph{current} one. A closure saves an environment but no store, and \code{closV} has no field for one.

\begin{keypoint}[Environment vs.\ store]
The environment follows the program text: it is saved in closures and restored when they are called, so a name means what it meant where it was written. The store follows time: it is passed from each step to the next in execution order and is never saved or restored, so a box's contents are whatever was last written, wherever the reading code was written.
\end{keypoint}
If closures saved the store too, \code{f} in \code{ex5} would read \code{a} in the store of the moment \code{f} was created, where it holds 1, and \code{ex5} would give 11, the wrong answer of Attempt 1. \Cref{sec:storeboxes} traces \code{ex5} and draws both tables.

\section{Boxes in the Store-Passing Interpreter}
\label{sec:storeboxes}

\subsection{The Box Clauses}

With the store in place, the four new clauses are short. \code{boxC} evaluates the initial contents, takes a fresh location, and adds a cell for it; \code{unboxC} evaluates the box and fetches its location:
\begin{lstlisting}
    [boxC (e)
          (type-case Result (eval-env env sto e)
            [res (v sto1)
                 (let ([l (new-loc)])
                   (res (boxV l) (override-store (cell l v) sto1)))])]
    [unboxC (e)
            (type-case Result (eval-env env sto e)
              [res (v1 sto1)
                   (let ([v2 (fetch (boxV-l v1) sto1)])
                     (res v2 sto1))])]
\end{lstlisting}
Note that \code{unboxC} fetches from \code{sto1}, the store \emph{after} evaluating \code{e}, in case evaluating the box expression itself wrote to it. \code{setboxC} evaluates the box and then the new contents, threading the store as \code{plusC} does, and adds a cell for the box's location on top of \code{sto2}:
\begin{lstlisting}
    [setboxC (e1 e2)
             (type-case Result (eval-env env sto e1)
               [res (v1 sto1)
                    (type-case Result (eval-env env sto1 e2)
                      [res (v2 sto2)
                           (res v1 (override-store (cell (boxV-l v1) v2) sto2))])])]
    [seqC (e1 e2)
          (type-case Result (eval-env env sto e1)
            [res (v1 sto1)
                 (eval-env env sto1 e2)])]
\end{lstlisting}
Two details:
\begin{itemize}
  \item \code{setboxC} returns \code{v1}, the box itself, not the new contents and not a void value (the language has none). This is a design choice; Racket's \code{set-box!} returns \code{(void)}.
  \item \code{seqC} discards \code{v1} but \emph{not} \code{sto1}: \code{e2} runs in the store that \code{e1} left behind. In Attempt~1 the \code{seqC} clause looked the same, but there was no store to pass on, so \code{e1} could only be evaluated and forgotten. Here sequencing finally has a purpose: \code{e1} is run for its effect on the store.
\end{itemize}

\subsection{Running the Examples}

Each in a fresh session, so that \code{new-loc} starts at 0:
\begin{lstlisting}[style=repl]
> (eval ex1)
- Result
(res (numV 4) '())
> (eval ex2)
- Result
(res (boxV 0) (list (cell 0 (numV 1)) (cell 0 (numV 0))))
\end{lstlisting}
\begin{lstlisting}[style=repl]
> (eval ex3)
- Result
(res (numV 0) (list (cell 0 (numV 0))))
> (eval ex4)
- Result
(res (numV 2) (list (cell 0 (numV 2)) (cell 0 (numV 1)) (cell 0 (numV 0))))
> (eval ex5)
- Result
(res (numV 12) (list (cell 0 (numV 2)) (cell 0 (numV 1))))
\end{lstlisting}
All five are now right. \code{ex1} allocates nothing, so its final store is empty. \code{ex2} is the result shown in lecture: the value is the box, location 0, and the store holds two cells for location 0, the write of 1 in front of the initial 0.

\subsection{Tracing \texttt{ex4}}

Write the store newest first, with \code{l:v} for \code{(cell l v)}. The outer \code{let} becomes an application of \code{(lambda b ...)} to \code{(box 0)}:
\begin{lstlisting}[style=repl]
(let b (box 0) (begin (begin (set-box! b 1)
                             (set-box! b (+ 1 (unbox b))))
                      (unbox b)))
  (lambda b ...)                          ---> closure, sto = []
  (box 0)       new location 0            sto = [0:0]
  body in env [b=(boxV 0)]
    (set-box! b 1)                        sto = [0:1, 0:0]
    (set-box! b (+ 1 (unbox b)))
      (unbox b): fetch 0 ---> 1
      (+ 1 1) ---> 2                      sto = [0:2, 0:1, 0:0]
    (unbox b): fetch 0 ---> 2
---> (numV 2) with sto = [0:2, 0:1, 0:0]
\end{lstlisting}
The environment never changes after the \code{let}: \code{b} is bound to \code{(boxV 0)} throughout. Every change happens in the store, and each \code{begin} hands the store from its first part to its second.

\subsection{Tracing \texttt{ex5}: Two Tables, Two Paths}

\begin{lstlisting}[style=repl]
(let a (box 1)
  (let f (lambda x (+ x (unbox a)))
    (begin (set-box! a 2) (f 10))))
  (box 1)       new location 0            sto = [0:1]
  env [a=(boxV 0)]
  (lambda x ...) ---> f = (closV [a=(boxV 0)] x (+ x (unbox a)))
  env [f=closure, a=(boxV 0)]
  (set-box! a 2)                          sto = [0:2, 0:1]
  (f 10)
    body env:   the closure's [a=(boxV 0)] + x=10
    body store: the current one,          sto = [0:2, 0:1]
    (unbox a): fetch 0 ---> 2
    (+ 10 2) ---> 12
---> (numV 12) with sto = [0:2, 0:1]
\end{lstlisting}
The closure saved the binding \code{a=(boxV 0)}: a location, not a number. The location has not changed since the closure was created, but the store has, and the call reads the current store (\Cref{fig:envstore}). This is how lexical scope and mutation coexist: the closure fixes \emph{which} box \code{a} is, and the store says what is in it \emph{now}.

\begin{figure}[htbp]
\centering
\begin{tikzpicture}[x=1cm, y=1cm]
  \node[cell, fill=tsk, minimum width=10mm] (cenv) at (0,0) {};
  \node[cell, fill=tsk, minimum width=7mm, right=0pt of cenv] (cx) {x};
  \node[cell, fill=tsk, minimum width=30mm, right=0pt of cx] (ce) {(+ x (unbox a))};
  \node[lbl, above=1mm of cx] {\code{f} = \code{closV}};
  \node[cell, fill=uspace, minimum width=22mm] (a) at (0,-1.6) {a=(boxV 0)};
  \fill (cenv.center) circle (0.6mm);
  \draw[ptr] (cenv.center) -- (a.north);
  \node[lbl, left=2mm of a, align=right, text=accent] {saved when the\\ \code{lambda} was evaluated};
  \node[cell, fill=kspace, minimum width=14mm] (x) at (0,-3.2) {x=10};
  \draw[ptr] (x) -- (a);
  \node[lbl, below=1mm of x, align=center] {environment of the body:\\ the closure's env + the parameter};
  \node[cell, fill=greybg, minimum width=12mm] (s2) at (5.6,-1.6) {0:2};
  \node[cell, fill=greybg, minimum width=12mm, right=5mm of s2] (s1) {0:1};
  \node[right=2mm of s1, lbl] (nil) {\code{[]}};
  \draw[ptr] (s2) -- (s1); \draw[ptr] (s1) -- (nil);
  \node[lbl, above=1mm of s2, xshift=8mm] {store \code{sto2}, passed in at the call};
  \draw[ptr, accent] (a.east) -- node[above, lbl, text=accent] {\code{fetch 0} stops here} (s2.west);
  \draw[decorate, decoration={brace, mirror, amplitude=4pt}] ($(s1.south west)+(0,-1mm)$) -- ($(nil.south east)+(0,-1mm)$) node[midway, below=2mm, lbl, align=center] {the store when\\ \code{f} was created};
\end{tikzpicture}
\caption{\code{ex5} while evaluating \code{(unbox a)} in \code{f}'s body. The environment (left) comes from the closure; the store (right) is the current one. The store at closure-creation time is still there as a suffix, but nothing reads it, since \code{fetch} stops at the newer cell.}
\label{fig:envstore}
\end{figure}

\begin{keypoint}[Where the environment and the store come from]
\begin{center}
\begin{tabular}{@{}p{0.3\linewidth}p{0.3\linewidth}p{0.3\linewidth}@{}}
\toprule
& environment & store \\
\midrule
to subexpressions & the same \code{env} to each & chained: \code{sto}, then \code{sto1}, \ldots \\
returned by \code{eval-env} & never & always, in the \code{res} \\
saved in a closure & yes & no \\
used for a function body & the closure's, extended & the current one, \code{sto2} \\
follows & the program text & the order of evaluation \\
\bottomrule
\end{tabular}
\end{center}
\end{keypoint}

\subsection{The Store Only Grows}

Nothing ever removes a cell. \code{set-box!} conses a new cell for an existing location, and the old one stays behind it, unreachable by \code{fetch}: after \code{ex4} the store holds three cells for location 0, of which only the first matters. Boxes that no part of the program can reach any more also keep their cells. A real implementation would overwrite a location in place and reclaim the memory of unreachable boxes with a garbage collector~\cite[Chapter~11]{PLAI}. For explaining what mutation \emph{means}, a store that only grows is simpler: every past state is still visible in the trace, and the store after any step is the store before it with some cells added on the front.

\section{The Meta-Circular Interpreter}
\label{sec:metacirc}

\subsection{Boxes Implemented by Boxes}

The store-passing interpreter needed a new type, a new argument and result for \code{eval-env}, and a rewrite of every clause. The lecture's second file starts again from the starting file and takes a shortcut: the object language's boxes are implemented with Racket's boxes.

\begin{definition}[Meta-circular interpreter]
An interpreter is \emph{meta-circular} when it implements a feature of the language it interprets by the same feature of the language it is written in. SICP uses the word for an evaluator written in the language that it evaluates~\cite[Section~4.1]{SICP}; here it applies feature by feature: first to boxes, then to functions.
\end{definition}

A box value holds a Racket box of values, \code{(boxof Value)}. \code{eval-env} goes back to the signature of \code{lecture5e.rkt}, \code{(env e) : Value}, with no store, and every old clause is unchanged:
\begin{lstlisting}
(define-type Value
  [numV (n : number)]
  [closV (env : Env) (x : symbol) (e : Expr)]
  [boxV (b : (boxof Value))])

    [boxC (e)
          (boxV (box (eval-env env e)))]
    [unboxC (e)
            (unbox (boxV-b (eval-env env e)))]
    [setboxC (e1 e2)
             (let ([v1 (eval-env env e1)]
                   [v2 (eval-env env e2)])
               (begin
                 (set-box! (boxV-b v1) v2)
                 v1))]
    [seqC (e1 e2)
          (begin
            (eval-env env e1)
            (eval-env env e2))]
\end{lstlisting}
Each clause does in Racket what the construct does in the object language. \code{(box e)} makes a Racket box holding the value of \code{e} and wraps it in \code{boxV}. \code{unbox} unwraps and reads it. \code{set-box!} evaluates the box and the new contents and mutates the Racket box; like the store-passing version, it returns the box \code{v1}. \code{seqC} is Racket's \code{begin}.

\begin{remark}
Once all eleven variants have clauses, the starting file's \code{else} clause can never be reached, and \code{plai-typed} rejects the program: ``the else branch of this type-case is unreachable''. Deleting the \code{else} fixes it.
\end{remark}

Compare this \code{setboxC} with Attempt~1's. There, \code{boxV} contained a \tyname{Value} directly, and \code{set-box!} could only build a new \code{boxV}. Here \code{boxV} contains a Racket box, a reference to a mutable cell. Copying the \code{boxV} into an environment or a closure copies the reference, not the cell, so every copy sees the same contents. This is the indirection of \Cref{sec:stateboxes}, supplied by Racket. The results agree with the store-passing interpreter on all five examples:
\begin{lstlisting}[style=repl]
> (eval ex1)
- Value
(numV 4)
> (eval ex2)
- Value
(boxV (box (numV 1)))
> (eval ex3)
- Value
(numV 0)
> (eval ex4)
- Value
(numV 2)
> (eval ex5)
- Value
(numV 12)
\end{lstlisting}
In \code{ex5}, the closure for \code{f} saves an environment in which \code{a} is bound to a \code{boxV} holding a Racket box; \code{(set-box! a 2)} mutates that Racket box, and the call reads it.

\subsection{Functions Implemented by Functions}

The lecture then applies the same idea to functions. A function value no longer records its parameter, body, and environment; it simply \emph{is} a Racket function from values to values. The constructor keeps the name \code{closV}:
\begin{lstlisting}
(define-type Value
  [numV (n : number)]
  [closV (f : (Value -> Value))]
  [boxV (b : (boxof Value))])

    [lambdaC (x e)
             (closV (lambda (arg-value)
                      (eval-env (extend-env (bind x arg-value) env) e)))]
    [appC (e1 e2)
          (let [(v1 (eval-env env e1))
                (v2 (eval-env env e2))]
            ((closV-f v1) v2))]
\end{lstlisting}
\code{lambdaC} builds a Racket \code{lambda} that, given the argument's value, evaluates the body in \code{env} extended with \code{x}. \code{appC} evaluates the function and the argument and calls the Racket function. The three accessors \code{closV-env}, \code{closV-x}, and \code{closV-e} of \code{lecture5e.rkt} are gone.

The \code{env} inside the Racket \code{lambda} is the environment of the \code{lambdaC} clause, the one where the object-language \code{lambda} was evaluated. Racket's own closure saves it, so the object language gets lexical scope without any code for it:
\begin{lstlisting}[style=repl]
> (eval (parse '((lambda x (+ 2 x)) 1)))
- Value
(numV 3)
> (eval (parse '(let x 2 (let addx (lambda y (+ y x)) (let x 6 (addx 3))))))
- Value
(numV 5)
> (eval (parse '(lambda x x)))
- Value
(closV #<procedure>)
\end{lstlisting}
The \code{addx} example of \Cref{sec:closures} gives 5, the lexical answer, and \code{ex5} still gives 12. A function value now prints as an opaque Racket procedure: its parameter, body, and saved environment can no longer be inspected. In week~5 it took two files, \code{lecture5d.rkt} and \code{lecture5e.rkt}, to get closures right; here the right answer comes for free, and for the same reason as the boxes, it explains nothing: the interpreter has lexical scope because Racket does.

\subsection{Side by Side}

\begin{center}
\small
\begin{tabular}{@{}p{0.2\linewidth}p{0.36\linewidth}p{0.36\linewidth}@{}}
\toprule
& store-passing & meta-circular \\
\midrule
box value & \code{(boxV l)}, a location & \code{(boxV b)}, a Racket box \\
function value & \code{(closV env x e)} & \code{(closV f)}, a Racket function \\
contents live in & the store, an explicit list of cells & Racket's heap \\
\code{eval-env} & \code{(env sto e) : Result} & \code{(env e) : Value} \\
\code{set-box!} & adds a cell to the store it returns & Racket's \code{set-box!} \\
\code{begin} & passes \code{sto1} from \code{e1} to \code{e2} & Racket's \code{begin} \\
order of evaluation & forced by the store: \code{e2} needs \code{sto1} & whatever Racket does \\
Racket mutation used & only in \code{new-loc} & for every box operation \\
old clauses & all rewritten to thread the store & unchanged, except \code{lambdaC} and \code{appC} once functions are meta-circular \\
\bottomrule
\end{tabular}
\end{center}

\subsection{What Each Version Explains}

The meta-circular interpreter is much shorter, and it is correct, but it explains mutation \emph{by} mutation. Someone who does not already know what \code{set-box!} does learns nothing from \code{[setboxC ... (set-box! (boxV-b v1) v2) ...]}. Two things in particular are inherited from Racket without appearing in the code:
\begin{itemize}
  \item \textbf{The store.} Where the contents of a box live, and what happens to them over time, is Racket's business. The store-passing interpreter writes the store down as a value that can be printed, as in the \code{(res ...)} results above.
  \item \textbf{The order of evaluation.} In \code{setboxC}, \code{e1} is evaluated before \code{e2} only because Racket's \code{let} evaluates its right-hand sides left to right; in \code{plusC}, only because Racket evaluates the arguments of \code{+} left to right. In the store-passing interpreter the order is forced by data dependencies: \code{e2} cannot be evaluated until \code{sto1} exists, and \code{sto1} is produced by evaluating \code{e1}.
\end{itemize}
Once there is mutation, order matters. Consider (not a lecture example) the program
\begin{lstlisting}
(let b (box 1) (+ (unbox b) (begin (set-box! b 10) (unbox b))))
\end{lstlisting}
Both interpreters give \code{(numV 11)}: the left operand reads 1 before the right operand writes 10. Evaluating the right operand first would give 20. Without mutation the two orders could never be told apart; with it, the order is part of the meaning of \code{+}, and the meta-circular interpreter specifies it only by reference to Racket. Reynolds pointed out this kind of dependence: a definitional interpreter can silently inherit the evaluation order of the language it is written in~\cite{Reynolds}.

\begin{keypoint}[What an interpreter explains]
An interpreter explains a feature only if it implements it with something simpler. Week~5 implemented environments with Racket lists and \code{plusC} with Racket's \code{+}; that was fine, because lists and addition were not what was being studied, and scope was implemented explicitly. This week the object of study is mutation, so implementing it with Racket's mutation is circular; implementing functions with Racket's functions likewise hides how closures capture their environment. The store-passing interpreter explains mutation with immutable lists and function arguments and results, and in doing so also makes evaluation order explicit.
\end{keypoint}

There is one remaining circle in the store-passing interpreter: \code{new-loc} uses a Racket box. It is not essential: \code{new-loc} only has to produce a location not yet in the store, and a counter threaded alongside the store would do. It is the same kind of shortcut as \code{gensym} in week~4 (\Cref{sec:capavoid}), a source of fresh names rather than the subject of the lecture.

